% GENERATED FILE. Edit canonical lessons, then run npm run generate:books.
% canonical-source-hash: fnv1a64:382eaef229fb5eb7
% canonical-lessons: SA-C317-dutah, SA-C317-ajna, SA-C317-adisati, SA-C317-sandhih, SA-C317-samantah, SA-R317-review-310-317

\chapter{The Envoy}
\label{ch:sa-the-envoy}
\hlchaptermodality{317}

\begin{chapteropening}
\textbf{By the end of this chapter:} \emph{I can tell how a king sent orders and made peace.}

Complete the last lesson of chapter 317: I can tell how a king sent orders and made peace.
\end{chapteropening}

\section[\texorpdfstring{\sk{दूतः} (dūtaḥ)}{दूतः (dūtaḥ)}]{\sk{दूतः} (dūtaḥ) — an envoy; a messenger}

\label{lesson:SA-C317-dutah}



\begin{quote}
Before the new one: say the Sanskrit for a kingdom, then the Sanskrit for governs.
\end{quote}

\subsection*{You'll want to know: \sk{दूतः}}
\textbf{\sk{दूतः}} — \emph{dūtaḥ} — \textquotedblleft{}an envoy; a messenger\textquotedblright{}.

\begin{quote}
\sk{शत्रोः} \sk{दूतः} \sk{सभाम्} \sk{आगच्छति।}
\end{quote}

\begin{quote}
\emph{śatroḥ dūtaḥ sabhām āgacchati.}
\end{quote}

\begin{quote}
The enemy's envoy comes to the court.
\end{quote}

A masculine noun.

\begin{cousinweb}
\textbf{\sk{दूतः}} (\emph{dūtaḥ}) MW thinks it probably comes from the root √du, \textquotedblleft{}to go\textquotedblright{}, and compares \emph{dūra}, \textquotedblleft{}far\textquotedblright{}; MW gives \emph{dūta} as a messenger, envoy, ambassador (MW).
\end{cousinweb}

\begin{grammarlens}[title={What you've built}]
One more word for telling what happened.
\end{grammarlens}

\subsection*{Your turn}
\begin{itemize}
\raggedright
  \item \emph{Say it:} \emph{dūtaḥ}
  \item \emph{Say it:} \emph{dūtaḥ}, once more
  \item \emph{Say it:} \emph{śatroḥ dūtaḥ sabhām āgacchati.} — The enemy's envoy comes to the court.
  \item \emph{Recall:} say the Sanskrit for a kingdom, then the Sanskrit for governs, then say \emph{dūtaḥ} again
\end{itemize}

\begin{tcolorbox}[breakable,colback=teal!4,colframe=teal!35!black,title={Before you move on}]
What does \sk{दूतः} mean? (\textquotedblleft{}An envoy; a messenger\textquotedblright{}.) Say it once more.
\end{tcolorbox}
\section[\texorpdfstring{\sk{आज्ञा} (ājñā)}{आज्ञा (ājñā)}]{\sk{आज्ञा} (ājñā) — an order; a command}

\label{lesson:SA-C317-ajna}



\begin{quote}
Before the new one: say the Sanskrit for governs, then the Sanskrit for an envoy.
\end{quote}

\subsection*{You'll want to know: \sk{आज्ञा}}
\textbf{\sk{आज्ञा}} — \emph{ājñā} — \textquotedblleft{}an order; a command\textquotedblright{}.

\begin{quote}
\sk{सैनिकाः} \sk{सेनापतेः} \sk{आज्ञां} \sk{पालयन्ति।}
\end{quote}

\begin{quote}
\emph{sainikāḥ senāpateḥ ājñāṃ pālayanti.}
\end{quote}

\begin{quote}
The soldiers obey the general's order.
\end{quote}

A feminine noun.

\begin{cousinweb}
\textbf{\sk{आज्ञा}} (\emph{ājñā}) is built from \emph{ā} and the root √jñā, \textquotedblleft{}to know\textquotedblright{}; MW gives \emph{ājñā} as an order or command (MW).
\end{cousinweb}

\begin{grammarlens}[title={What you've built}]
One more word for telling what happened.
\end{grammarlens}

\subsection*{Your turn}
\begin{itemize}
\raggedright
  \item \emph{Say it:} \emph{ājñā}
  \item \emph{Say it:} \emph{ājñā}, once more
  \item \emph{Say it:} \emph{sainikāḥ senāpateḥ ājñāṃ pālayanti.} — The soldiers obey the general's order.
  \item \emph{Recall:} say the Sanskrit for governs, then the Sanskrit for an envoy, then say \emph{ājñā} again
\end{itemize}

\begin{tcolorbox}[breakable,colback=teal!4,colframe=teal!35!black,title={Before you move on}]
What does \sk{आज्ञा} mean? (\textquotedblleft{}An order; a command\textquotedblright{}.) Say it once more.
\end{tcolorbox}
\section[\texorpdfstring{\sk{आदिशति} (ādiśati)}{आदिशति (ādiśati)}]{\sk{आदिशति} (ādiśati) — commands; gives orders}

\label{lesson:SA-C317-adisati}



\begin{quote}
Before the new one: say the Sanskrit for an envoy, then the Sanskrit for an order.
\end{quote}

\subsection*{You'll want to know: \sk{आदिशति}}
\textbf{\sk{आदिशति}} — \emph{ādiśati} — \textquotedblleft{}commands; gives orders\textquotedblright{}.

\begin{quote}
\sk{राजा} \sk{सेनापतिम्} \sk{आदिशति।}
\end{quote}

\begin{quote}
\emph{rājā senāpatim ādiśati.}
\end{quote}

\begin{quote}
The king gives orders to the general.
\end{quote}

Past (he …): \sk{सः} \sk{सेनापतिम्} \sk{आदिष्टवान्}, \emph{saḥ senāpatim ādiṣṭavān}.

\begin{cousinweb}
\textbf{\sk{आदिशति}} (\emph{ādiśati}) is built from \emph{ā} and the root √diś, \textquotedblleft{}to point, show\textquotedblright{}, as in \emph{diśā}; MW gives \emph{ā-diś} as to point out, assign, and to order, direct, command (MW).
\end{cousinweb}

\begin{grammarlens}[title={What you've built}]
One more word for telling what happened.
\end{grammarlens}

\subsection*{Your turn}
\begin{itemize}
\raggedright
  \item \emph{Say it:} \emph{ādiśati}
  \item \emph{Say it:} \emph{ādiśati}, once more
  \item \emph{Say it:} \emph{rājā senāpatim ādiśati.} — The king gives orders to the general.
  \item \emph{Recall:} say the Sanskrit for an envoy, then the Sanskrit for an order, then say \emph{ādiśati} again
\end{itemize}

\begin{tcolorbox}[breakable,colback=teal!4,colframe=teal!35!black,title={Before you move on}]
What does \sk{आदिशति} mean? (\textquotedblleft{}Commands; gives orders\textquotedblright{}.) Say it once more.
\end{tcolorbox}
\section[\texorpdfstring{\sk{सन्धिः} (sandhiḥ)}{सन्धिः (sandhiḥ)}]{\sk{सन्धिः} (sandhiḥ) — a peace treaty; an alliance}

\label{lesson:SA-C317-sandhih}



\begin{quote}
Before the new one: say the Sanskrit for an order, then the Sanskrit for commands.
\end{quote}

\subsection*{You'll want to know: \sk{सन्धिः}}
\textbf{\sk{सन्धिः}} — \emph{sandhiḥ} — \textquotedblleft{}a peace treaty; an alliance\textquotedblright{}.

\begin{quote}
\sk{अन्ते} \sk{शत्रुः} \sk{सन्धिं} \sk{याचते।}
\end{quote}

\begin{quote}
\emph{ante śatruḥ sandhiṃ yācate.}
\end{quote}

\begin{quote}
In the end the enemy asks for a peace treaty.
\end{quote}

A masculine noun.

\begin{cousinweb}
\textbf{\sk{सन्धिः}} (\emph{sandhiḥ}) is built from \emph{sam}, \textquotedblleft{}together\textquotedblright{}, and the root √dhā, \textquotedblleft{}to put\textquotedblright{}: a putting together, a junction, and MW also gives \emph{saṃdhi} as an alliance or peace; in grammar it is the joining of sounds (MW).
\end{cousinweb}

\begin{grammarlens}[title={What you've built}]
One more word for telling what happened.
\end{grammarlens}

\subsection*{Your turn}
\begin{itemize}
\raggedright
  \item \emph{Say it:} \emph{sandhiḥ}
  \item \emph{Say it:} \emph{sandhiḥ}, once more
  \item \emph{Say it:} \emph{ante śatruḥ sandhiṃ yācate.} — In the end the enemy asks for a peace treaty.
  \item \emph{Recall:} say the Sanskrit for an order, then the Sanskrit for commands, then say \emph{sandhiḥ} again
\end{itemize}

\begin{tcolorbox}[breakable,colback=teal!4,colframe=teal!35!black,title={Before you move on}]
What does \sk{सन्धिः} mean? (\textquotedblleft{}A peace treaty; an alliance\textquotedblright{}.) Say it once more.
\end{tcolorbox}
\section[\texorpdfstring{\sk{सामन्तः} (sāmantaḥ)}{सामन्तः (sāmantaḥ)}]{\sk{सामन्तः} (sāmantaḥ) — a vassal; a feudal chief}

\label{lesson:SA-C317-samantah}



\begin{quote}
Before the new one: say the Sanskrit for commands, then the Sanskrit for a peace treaty.
\end{quote}

\subsection*{You'll want to know: \sk{सामन्तः}}
\textbf{\sk{सामन्तः}} — \emph{sāmantaḥ} — \textquotedblleft{}a vassal; a feudal chief\textquotedblright{}.

\begin{quote}
\sk{सामन्ताः} \sk{सभायाम्} \sk{उपविशन्ति।}
\end{quote}

\begin{quote}
\emph{sāmantāḥ sabhāyām upaviśanti.}
\end{quote}

\begin{quote}
The vassals sit in the assembly.
\end{quote}

A masculine noun.

\begin{cousinweb}
\textbf{\sk{सामन्तः}} (\emph{sāmantaḥ}) is built on \emph{sam-anta}, \textquotedblleft{}bordering\textquotedblright{}, with the lengthened first vowel: MW gives \emph{sāmanta} as being on all sides, a neighbour, and a vassal or feudal prince (MW).
\end{cousinweb}

\begin{grammarlens}[title={What you've built}]
One more word for telling what happened.
\end{grammarlens}

\subsection*{Your turn}
\begin{itemize}
\raggedright
  \item \emph{Say it:} \emph{sāmantaḥ}
  \item \emph{Say it:} \emph{sāmantaḥ}, once more
  \item \emph{Say it:} \emph{sāmantāḥ sabhāyām upaviśanti.} — The vassals sit in the assembly.
  \item \emph{Recall:} say the Sanskrit for commands, then the Sanskrit for a peace treaty, then say \emph{sāmantaḥ} again
\end{itemize}

\begin{tcolorbox}[breakable,colback=teal!4,colframe=teal!35!black,title={Before you move on}]
What does \sk{सामन्तः} mean? (\textquotedblleft{}A vassal; a feudal chief\textquotedblright{}.) Say it once more.
\end{tcolorbox}
\section[(dialogue)]{Looking back: from What Happened Today to The Envoy}

\label{lesson:SA-R317-review-310-317}



\begin{quote}
Say the Sanskrit for an event and befalls.
\end{quote}

\subsection*{Your turn}
Say these aloud:

\begin{itemize}
\raggedright
  \item an event, befalls, a detailed account, publishes, a detail
  \item a rumour, dread, stirs up, doubts, a misunderstanding
  \item a spokesman, an allegation, denies, a lie, doubtful
  \item a disturbance, a quarrel, unrest, surrounds, quells
  \item valour, steadfast, fame, honour, congratulates
  \item an explosion, an echo, unexpected, is alarmed, hurries
  \item an emperor, a throne, a lord, a kingdom, governs
  \item an envoy, an order, commands, a peace treaty, a vassal
\end{itemize}

\begin{tcolorbox}[breakable,colback=teal!4,colframe=teal!35!black,title={Before you move on}]
An event? (\textbf{\sk{घटना}}, \emph{ghaṭanā}.) Befalls? (\textbf{\sk{आपतति}}, \emph{āpatati}.) A detailed account? (\textbf{\sk{विवरणम्}}, \emph{vivaraṇam}.) Publishes? (\textbf{\sk{प्रकाशयति}}, \emph{prakāśayati}.) A detail? (\textbf{\sk{विशेषः}}, \emph{viśeṣaḥ}.) A rumour? (\textbf{\sk{किंवदन्ती}}, \emph{kiṃvadantī}.) Dread? (\textbf{\sk{आतंकः}}, \emph{ātaṅkaḥ}.) Stirs up? (\textbf{\sk{उत्तेजयति}}, \emph{uttejayati}.) Doubts? (\textbf{\sk{शंकते}}, \emph{śaṅkate}.) A misunderstanding? (\textbf{\sk{भ्रान्तिः}}, \emph{bhrāntiḥ}.) A spokesman? (\textbf{\sk{प्रवक्ता}}, \emph{pravaktā}.) An allegation? (\textbf{\sk{आरोपः}}, \emph{āropaḥ}.) Denies? (\textbf{\sk{निराकरोति}}, \emph{nirākaroti}.) A lie? (\textbf{\sk{अनृतम्}}, \emph{anṛtam}.) Doubtful? (\textbf{\sk{सन्दिग्धः}}, \emph{sandigdhaḥ}.) A disturbance? (\textbf{\sk{उपद्रवः}}, \emph{upadravaḥ}.) A quarrel? (\textbf{\sk{कलहः}}, \emph{kalahaḥ}.) Unrest? (\textbf{\sk{क्षोभः}}, \emph{kṣobhaḥ}.) Surrounds? (\textbf{\sk{वेष्टयति}}, \emph{veṣṭayati}.) Quells? (\textbf{\sk{शमयति}}, \emph{śamayati}.) Valour? (\textbf{\sk{शौर्यम्}}, \emph{śauryam}.) Steadfast? (\textbf{\sk{धीरः}}, \emph{dhīraḥ}.) Fame? (\textbf{\sk{कीर्तिः}}, \emph{kīrtiḥ}.) Honour? (\textbf{\sk{गौरवम्}}, \emph{gauravam}.) Congratulates? (\textbf{\sk{अभिनन्दति}}, \emph{abhinandati}.) An explosion? (\textbf{\sk{विस्फोटः}}, \emph{visphoṭaḥ}.) An echo? (\textbf{\sk{प्रतिध्वनिः}}, \emph{pratidhvaniḥ}.) Unexpected? (\textbf{\sk{आकस्मिकः}}, \emph{ākasmikaḥ}.) Is alarmed? (\textbf{\sk{उद्विजते}}, \emph{udvijate}.) Hurries? (\textbf{\sk{त्वरते}}, \emph{tvarate}.) An emperor? (\textbf{\sk{सम्राट्}}, \emph{samrāṭ}.) A throne? (\textbf{\sk{सिंहासनम्}}, \emph{siṃhāsanam}.) A lord? (\textbf{\sk{ईश्वरः}}, \emph{īśvaraḥ}.) A kingdom? (\textbf{\sk{राज्यम्}}, \emph{rājyam}.) Governs? (\textbf{\sk{शास्ति}}, \emph{śāsti}.) An envoy? (\textbf{\sk{दूतः}}, \emph{dūtaḥ}.) An order? (\textbf{\sk{आज्ञा}}, \emph{ājñā}.) Commands? (\textbf{\sk{आदिशति}}, \emph{ādiśati}.) A peace treaty? (\textbf{\sk{सन्धिः}}, \emph{sandhiḥ}.) A vassal? (\textbf{\sk{सामन्तः}}, \emph{sāmantaḥ}.)
\end{tcolorbox}
